\section{Method}
\label{sec:method}

\subsection{Architecture}

The detector is a frozen MERT-v1-95M backbone with a small trainable head.
Audio is resampled to \SI{24}{\kilo\hertz} mono and segmented into fixed clips;
MERT returns thirteen hidden states per clip at \SI{74.8}{\hertz}, each of
dimension 768. The head aggregates these in three stages.

\paragraph{Per-layer normalisation.}
The thirteen states already share a shape, but differ substantially in scale.
Without normalisation the layer-attention softmax tracks magnitude rather than
informativeness, so each state is passed through its own layer normalisation
before weighting.

\paragraph{Layer attention.}
A learned scalar per layer, softmax-normalised, produces a weighted sum over
depth. Learning one global profile rather than a per-example weighting keeps the
weights readable as an explanation: they state which depths the detector relies
on in general.

\paragraph{Temporal and clip pooling.}
Additive attention pools frames within a clip, and a second attention stage
pools clips within a song. The second stage is necessary rather than
decorative. SONICS songs run \SIrange{32}{240}{\second}, MERT was pre-trained on
\SI{5}{\second} clips, and its attention is $O(T^2)$: at \SI{74.8}{\hertz} a
\SI{120}{\second} song is roughly 9000 frames, whose attention matrix alone
exceeds \SI{1.9}{\giga\byte} per layer in half precision. Audio must therefore
be segmented, and without a song-level stage the model emits per-clip decisions
while the task, and the published baselines, are per-song.

The backbone contributes \SI{94.4}{\mega\nothing} frozen parameters and the head
\SI{0.62}{\mega\nothing} trainable ones.

\subsection{Classification heads and calibration}

The primary output is a binary real/synthetic logit; an auxiliary head predicts
the four-way SONICS taxonomy and is down-weighted so it cannot dominate the
objective that is reported. Probabilities are calibrated by temperature scaling
fitted on validation, never on the evaluation set.

Choosing the operating point required care. When validation separates the
classes completely---which it does here---the conventional choice of the
$(1-\alpha)$ quantile of scores on real music places the threshold immediately
against the real class. That threshold is valid on validation and catastrophic
elsewhere: applied to a shifted distribution it labels everything synthetic. We
instead place the threshold at the midpoint of the separating margin, which
scores identically on validation but is maximally distant from both classes and
therefore far more likely to survive a shift. Section~\ref{sec:results}
quantifies the difference: the quantile threshold yields FPR 1.00 on the
cross-corpus evaluation, the midpoint threshold FPR 0.00 on the cross-generator
one.

\subsection{Explanation layer}

Three levels are produced for every explained song.

\textbf{Level 1} reports the learned attention weights over layers, frames and
clips, which the architecture yields at no cost.

\textbf{Level 2} localises evidence in the signal. MERT consumes raw waveform
through a one-dimensional convolutional stack whose output is
$512~\text{channels} \times \text{time}$; no time--frequency map exists anywhere
in the forward graph. Gradient-based attribution taken there localises
\emph{when}, not \emph{at which frequency}, and we report it as such (Level 2a).
For frequency evidence we occlude tiles in the short-time Fourier domain and
resynthesise, which measures what the model does when a band is actually absent
(Level 2b); this is expensive and is applied to an audit subset.

\textbf{Level 3} computes exact Shapley values over the thirteen layers, treating
each hidden state as a player and the logit under a renormalised subset as the
coalition value. With thirteen layers the $2^{13}$ coalitions are affordable
because the frozen backbone runs once and only the small head is re-evaluated.

Levels 1 and 3 answer different questions---what the model weights, versus what
moved a particular decision---and we report their disagreement rather than
averaging it away.

\subsection{Explanation faithfulness}

No explanation is reported without being scored. We mask the regions an
explanation marks important and measure the movement in the logit against
masking randomly chosen regions of the same extent. The quantity reported is the
\emph{gap} against that random baseline, not the raw deletion curve: masking any
audio degrades a detector somewhat, so a steep curve alone establishes nothing.
